% =====================================================================
%  1bat-ud1-8.tex  —  HORA 8: Consolidació d'inequacions (casos difícils i errors típics)
% =====================================================================
\horatitol{Sessió 8 — Consolidació d'inequacions: casos difícils i errors típics}%
{Reforç abans dels contextos amb diverses condicions: caçar els errors habituals i tractar els casos amb truc.}

\subsection{Idea clau}

Avui no hi ha matèria nova: consolidem la resolució d'inequacions caçant els
\textbf{errors típics} i tractant els \textbf{casos amb truc}. L'objectiu és que els
passos crítics —girar el signe, treure parèntesis, operar amb fraccions— surtin
sense vacil·lar.

\begin{idea}
\textbf{Els errors que cacem avui:}
\begin{itemize}[leftmargin=1.4em, itemsep=2pt]
  \item \textbf{No girar el signe} en dividir o multiplicar per un negatiu.
  \item \textbf{Treure malament un parèntesi} (oblidar de canviar els signes de dins).
  \item \textbf{Perdre un denominador} en operar amb fraccions.
\end{itemize}
\textbf{Casos amb truc:} a vegades la $x$ \emph{desapareix} en simplificar i queda
una afirmació sobre números. Aleshores:
\begin{itemize}[leftmargin=1.4em, itemsep=2pt]
  \item Si queda \textbf{sempre certa} (com $1>-3$), la solució és \textbf{tots els
        reals}, $\mathbb{R}$.
  \item Si queda \textbf{mai certa} (com $0>5$), \textbf{no hi ha solució}.
\end{itemize}
\end{idea}

\subsection{Exemples}

\begin{exemple}[caça l'error: el signe]
Algú resol $-3x<9$ així: «$x<\dfrac{9}{-3}=-3$, per tant $x<-3$». \textbf{És
incorrecte:} en dividir per $-3$ (negatiu) calia \textbf{girar} el símbol. El correcte:
\[
-3x<9 \;\Rightarrow\; x>\frac{9}{-3} \;\Rightarrow\; x>-3.
\]
\textbf{Comprovació:} $x=0$ (que compleix $x>-3$) dóna $-3\cdot 0=0<9$, cert.
\end{exemple}

\begin{exemple}[un cas amb truc: sempre cert]
Resolem $2x+1>2x-3$. Restem $2x$ als dos costats:
\[
1>-3.
\]
La $x$ ha desaparegut i queda $1>-3$, que és \textbf{sempre cert}, sigui quina sigui
la $x$. Per tant, la solució és \textbf{tots els reals}: $x\in\mathbb{R}$, interval
$(-\infty,+\infty)$.
\end{exemple}

\subsection{Exercicis resolts}

\begin{exresolt}[detectar i corregir un error]
Aquesta resolució té un error. Troba'l i corregeix-la:
\[
2(x-3)\ge 4x \;\Rightarrow\; 2x-3\ge 4x \;\Rightarrow\; -2x\ge 3 \;\Rightarrow\; x\ge -\tfrac{3}{2}.
\]
\solucio
Hi ha \textbf{dos} errors. El primer, en treure el parèntesi: $2(x-3)=2x-6$, no
$2x-3$. El segon, al final: en dividir per $-2$ (negatiu) cal \textbf{girar} el
símbol. Resolució correcta:
\[
2(x-3)\ge 4x \;\Rightarrow\; 2x-6\ge 4x \;\Rightarrow\; -2x\ge 6 \;\Rightarrow\; x\le -3.
\]
Interval $(-\infty,-3]$.
\resposta{Errors: $2(x-3)=2x-6$ i no girar el signe. Solució correcta: $x\le -3$.}
\end{exresolt}

\begin{exresolt}[cas sense solució]
Resol $3x+5<3x-2$.
\solucio
Restem $3x$ als dos costats:
\[
5<-2.
\]
La $x$ desapareix i queda $5<-2$, que és \textbf{fals} sempre. Per tant, la inequació
\textbf{no té solució}: cap valor de $x$ la compleix.
\resposta{No té solució (queda $5<-2$, sempre fals).}
\end{exresolt}

\subsection{Exercicis per fer}

\noindent\textit{Itinerari: 1–4 són la base; 5 i 6 són ampliació (fraccions i casos
amb truc).}

\begin{exercicis}
\begin{exlist}
  \item \textbf{(Caça l'error)} Digues si cada resolució és correcta i, si no, corregeix-la:
        \begin{enumerate}[label=\alph*), leftmargin=1.6em, topsep=2pt]
          \item «$-2x<6 \Rightarrow x<-3$»
          \item «$3(x-1)\le x \Rightarrow 3x-1\le x \Rightarrow 2x\le 1 \Rightarrow x\le \tfrac{1}{2}$»
        \end{enumerate}
  \item Resol i dóna l'interval: $5x-3\le 2x+9$.
  \item Resol i dóna l'interval: $-x+7>2x-2$ \quad (compte amb el signe).
  \item Resol i dóna l'interval: $2(x+1)-3\ge x$.
  \item \textbf{(Fraccions)} Resol: $\dfrac{x}{2}+1\le \dfrac{x-1}{3}$.
  \item \textbf{(Casos amb truc)} Resol i digues si la solució és «tots els reals» o
        «cap»:
        \begin{enumerate}[label=\alph*), leftmargin=1.6em, topsep=2pt]
          \item $4x+2>4x-1$
          \item $x-5\ge x+1$
        \end{enumerate}
\end{exlist}
\end{exercicis}

% ---------------------------------------------------------------------
\begin{idea}
\textbf{Full d'autoavaluació.} Marca amb sinceritat com et trobes, de cara a la prova:
\begin{center}
\renewcommand{\arraystretch}{1.4}
\begin{tabular}{p{8cm} c c c}
\toprule
\textbf{Sé fer\dots} & \textbf{Sol} & \textbf{Amb ajuda} & \textbf{Encara no} \\
\midrule
Resoldre una inequació senzilla (sense canvi de signe) & & & \\
Girar el símbol en dividir per un negatiu & & & \\
Treure parèntesis correctament & & & \\
Operar amb fraccions (treure denominadors) & & & \\
Escriure la solució en forma d'interval & & & \\
\bottomrule
\end{tabular}
\end{center}
\end{idea}

% ---------------------------------------------------------------------
\fullsolucions{Sessió 8}
\begin{solucions}
\begin{sollist}
  \item (a) \textbf{Incorrecta}: en dividir per $-2$ cal girar; $x>-3$. \quad
        (b) \textbf{Incorrecta}: $3(x-1)=3x-3$, no $3x-1$; correcte:
        $3x-3\le x \Rightarrow 2x\le 3 \Rightarrow x\le\tfrac{3}{2}$.
  \item $3x\le 12 \Rightarrow x\le 4$, interval $(-\infty,4]$.
  \item $-x-2x>-2-7 \Rightarrow -3x>-9$; girem: $x<3$, interval $(-\infty,3)$.
  \item $2x+2-3\ge x \Rightarrow 2x-1\ge x \Rightarrow x\ge 1$, interval $[1,+\infty)$.
  \item (Mult.\ per $6$) $3x+6\le 2(x-1) \Rightarrow 3x+6\le 2x-2 \Rightarrow x\le -8$,
        interval $(-\infty,-8]$.
  \item (a) Queda $2>-1$, sempre cert: solució \textbf{tots els reals}, $\mathbb{R}$. \quad
        (b) Queda $-5\ge 1$, sempre fals: \textbf{cap} solució.
\end{sollist}
\end{solucions}
